\documentclass{article}
\usepackage{amsmath,amssymb}
\usepackage{color, graphics}
\usepackage{graphicx}
\usepackage{algorithmic}
\usepackage{algorithm}
\usepackage{verbatim}
\usepackage{fullpage}
\usepackage{hyperref}
\usepackage{currfile}
\usepackage{esvect}
\usepackage[utf8]{inputenc}
\usepackage[english]{babel}
\begin{document}
%%%%%%%%%%%%%%%%%%%%%%%%%%%%%%%%%%%%%%%%%%%%%%%%%%%%%%%%%%%%%%%%%%%%

~
\vskip 4cm
\begin{center}
\Huge{EXCILL2: Extremal Combinatorics at Illinois}\\
\Large{March 16--18, 2013}\\
\Large{A conference hosted by the Department of Mathematics}\\
\Large{University of Illinois at Urbana-Champaign}\\
\vskip 1em
\url{http://www.math.uiuc.edu/excill2/}
\end{center}
\thispagestyle{empty}
\begin{center}
\vskip 2cm
\includegraphics[width=4cm]{logo}
\end{center}
\vfill
\hfill Organizers: J. Balogh, A. Kostochka, B. Lidick\'y, and D.B. West
\vskip 5pt
\noindent Supported by: \\
 IMA (Institute for Mathematics and its Applications)\\
 NSF grant (DMS-0745185) \\
 Department of Mathematics of the
University
of Illinois at Urbana-Champaign
\vskip 1pt
\hfill \today

\newpage


\centerline {\LARGE {\bf Conference Program}}

\def\programentry#1#2{#1  \emph{#2} \\}

\textbf{Saturday, Mar 16}
\vskip -9pt
\line(1,0){250}

\begin{tabular}{ll}
 8:50 & \programentry{Opening remarks}{}
9:00 & \programentry{Béla Bollobás,}{Bootstrap Percolation - Problems, Results and Conjectures}
10:00 & \programentry{coffee}{}
10:25 & \programentry{Benny Sudakov,}{ Counting and packing Hamilton cycles}
11:00 & \programentry{Paul N. Balister,}{Controllability, matchings, and the Karp-Sipser algorithm on random graphs}
11:35 & \programentry{Penny Haxell,}{Extremal hypergraphs for packing and covering}
12:10 & \programentry{Andrzej Ruciński,}{On Two Hamiltonian Problems For Triple Systems}
12:45 & \programentry{lunch}{}
2:15 & \programentry{Alexander Razborov,}{Flag Algebras: an Interim Report}
3:15 & \programentry{Sergey Norin,}{Finitely-forcible graphons}
3:50 & \programentry{conference photo}{}
3:51 & \programentry{coffee}{}
4:15 & \programentry{Dhruv Mubayi,}{Hypergraph Ramsey Problems}
4:50 & \programentry{Tao Jiang,}{Hypergraph Turán and Ramsey results on linear cycles}
5:25 & \programentry{Jacques Verstraëte,}{On independent sets in locally sparse hypergraphs}
6:00 & \programentry{break}{}
6:15 & \programentry{Pizza reception and Poster session}{}

\end{tabular}

\textbf{Sunday, Mar 17}
\vskip -9pt
\line(1,0){250}

\begin{tabular}{ll}
9:00 & \programentry{Xuding Zhu,}{Total weight choosability of graphs}
10:00 & \programentry{coffee}{}
10:25 & \programentry{William T. Trotter,}{Tree-width and Dimension}
11:00 & \programentry{Vladimir Nikiforov,}{Asymptotic Ramsey numbers of books}
11:35 & \programentry{Gábor N. Sárközy,}{Coverings by monochromatic pieces}
12:10 & \programentry{Hal Kierstead,}{Progress on directed versions of the Hajnal-Szemerédi Theorem}
12:45 & \programentry{lunch}{}
2:15 & \programentry{Wojciech  Samotij,}{Typical structure of graph homomorphisms}
2:50 & \programentry{David Galvin,}{Some extremal questions for coloring}
3:25 & \programentry{coffee}{}
4:00 & \programentry{Choongbum Lee,}{Long paths and cycles in random subgraphs of graphs with large minimum degree}
4:35 & \programentry{John Lenz,}{Spectral extremal problems for hypergraphs}
5:10 & \programentry{Guantao Chen,}{Reed-Frost transition on dense graphs}
5:45 & \programentry{end}{}
6:30 & \programentry{Dinner reception}{}

\end{tabular}

\textbf{Monday, Mar 18}
\vskip -9pt
\line(1,0){250}

\begin{tabular}{ll}
9:00 & \programentry{Zoltán Füredi,}{A discrete isodiametric result: the Erdős-Ko-Rado theorem for multisets}
9:35 & \programentry{Daniel Král',}{Quasirandomness of permutations}
10:10 & \programentry{coffee}{}
10:35 & \programentry{Allan Siu Lun Lo,}{Proof of the 1-factorization and Hamilton decomposition conjectures}
11:10 & \programentry{Richard Mycroft,}{Perfect matchings in dense hypergraphs}
11:45 & \programentry{László Babai,}{Symmetry vs. regularity}
12:45 & \programentry{lunch}{}

\end{tabular}

\thispagestyle{empty}

\newpage

\def\pen{\penalty10000}

%% NAME EMAIL AFFILIATION
\def\speaker#1#2#3{\hbox\bgroup\vrule
  \vbox\bgroup\hrule\pen\vskip0.5mm\hrule\pen\smallskip\begin{center}{\Large #1}\pen\vskip1pt%
  {\tt #2}\pen\smallskip%
  \write\spk{\string\obsline {#1}{}{\the\count0}}\pen
% \label{#1}
%  \write\spk{\string\obsline {#1}{}{#1}}\pen
\\ #3  
  }

%% NAME EMAIL AFFILIATION
\def\speakerposter#1#2#3{\hbox\bgroup\vrule
  \vbox\bgroup\hrule\pen\vskip0.5mm\hrule\pen\smallskip\begin{center}{\Large #1}\pen\vskip1pt%
  {\tt #2}\pen\smallskip%
%  \write\spk{\string\obsline {#1}{}{\the\count0}}\pen
% \label{#1}
%  \write\spk{\string\obsline {#1}{}{#1}}\pen
\\ #3  
  }

\def\titlebegin{\medskip}
\def\titleline#1{{\\\Large #1}\pen}
\def\titleend{\smallskip}
\def\title#1{\titlebegin\titleline{\Large #1}\titleend\write\obs{\string\obsline {#1}{}{\the\count0}}}
\def\titleposter#1{\titlebegin\titleline{\Large #1}\titleend}
%\def\title#1{\titlebegin\titleline{\Large #1}\titleend\write\obs{\string\obsline {#1}{}{\currfilename}}}

\def\source#1{{\\(#1)}\pen}

\def\endtitle{\end{center}\smallskip\hrule\egroup\vrule\egroup\pen\smallskip\bgroup\setcounter{Defnum}{0}}


%%%%% Some Macros - leave them as they are  %%%%%
%\def\speaker#1#2{%
%\begin{spktitle}\pen%
%\centerline{\Large \textbf{#1}}\pen%
%\vskip 1pt\pen%
%\centerline{\tt #2}\nopagebreak%
%\write\spk {\string\obsline {#1}{}{\the\count0}}%
%}                                                                          
%\def\title#1{%
%\begin{center}%
%{\Large\nopagebreak \textbf{#1}}%
%\end{center}%
%\write\obs {\string\obsline {#1}{}{\the\count0}}%
%}
%\def\author#1{\nopagebreak\smallskip\centerline{Presented paper by #1}\nopagebreak}                                                       
%\def\source#1{\nopagebreak\centerline{(#1)}\nopagebreak}
%\def\endtitle{\end{spktitle}\par\setcounter{Defnum}{0}\nopagebreak}

% Nejaka makra pro definovani definic a podobne
\newcounter{Defnum}
\newenvironment{definition}{\stepcounter{Defnum}\par\textbf{Definition \arabic{Defnum}.}} {}
\newenvironment{theorem}{\stepcounter{Defnum}\par\textbf{Theorem \arabic{Defnum}.}\it}{}
\newenvironment{lemma}{\stepcounter{Defnum}\par\textbf{Lemma \arabic{Defnum}.}\it}{}
\newenvironment{observation}{\stepcounter{Defnum}\par\textbf{Observation \arabic{Defnum}.}\it}{}
\newenvironment{corollary}{\stepcounter{Defnum}\par\textbf{Corollary \arabic{Defnum}.}\it}{}
\newenvironment{conjecture}{\stepcounter{Defnum}\par\textbf{Conjecture \arabic{Defnum}.}\it}{}
\newenvironment{proof}{\par\textbf{Proof}}{ \rule{1pt}{0pt}   \hfill $\square$ \\}
\newenvironment{proposition}{\stepcounter{Defnum}\par\textbf{Proposition \arabic{Defnum}.}\it}{}

%% Pro seznam participantu
\def\part#1#2#3#4{%
\leftline{\textbf{#1}}\penalty5000
\leftline{\hskip1cm #2}\penalty5000
\leftline{\hskip1cm #3}\penalty5000
\leftline{\hskip1cm\tt #4}\medskip}


% Dela sekce v jednotlivych kouskach sbornicku
\def\subsection#1{\par Not good idea. \smallskip}
\def\section#1{\par\medskip\centerline{\bf #1}\smallskip}

%% Pro formatovani odstavcu
\setlength{\parindent}{0pt}
\setlength{\parskip}{0.5 em}



%% Makro pro vkladani prispevku
\def\pripevek#1{
%	\label{#1.tex}
	\input{#1}\vskip 3em
}

% Makro pro uvedeni prispevku
% pouziti: \prispevek{nazev}{autori}

\def\prispevek #1#2{%
\clearpage
%  \par\penalty-200\bigskip\bigskip\bigskip
  \vbox{%
  \bgroup
    \parindent=0pt
    \spaceskip=.3333em
    \rightskip=3em plus 1fil
    \leftskip=3em plus 1fil
    \parfillskip=0em
    {\LARGE {\bf #1} \par}%Large
    \ifx @#2@\else 
    \penalty10000%
    \medskip
    {\Large {\em #2} \par}%large
    \fi
  \egroup
  }%
  \penalty10000%
  %\leavevmode
  \medskip
  \penalty10000%
  \bigskip
%  \doobsahu{#1}{#2}%
  \penalty10000%
}

% Zapise do pomocneho souboru fikane makro, 
% cimz zpusobi vytisteni prislusneho radku v obsahu.
% pouziti: \doobsahu{nazev}{autori}

\def\doobsahu #1#2{%
  \write\obs {\string\obsline {#1}{}{\the\count0}}%
  \write\spk {\string\obsline {#2}{}{\the\count0}}%
  \penalty 10000%
}

\def\prispevekvarianta #1#2{%
  \section*{#1}%
  \penalty10000%
  \medskip
  \centerline{\it #2}%
  \penalty10000%
  \medskip
  \bigskip
  \penalty10000%
% \doobsahu{#1}{#2}%
}

\def\participant#1#2#3#4{%
#1 #2 \url{#3} #4  \vskip -2pt
}

% Makro, ktere vyrobi jednu radku obsahu,
% je obsazeno v automaticky vygenerovanem souboru \jobname.obs
% pouziti: \obsline{nazev}{autori}{cislo stranky}

\def\obsline #1#2#3{%
\ifx @#2@\else #2: \fi
  #1 \leaders\hbox to 1em{\hfil.\hfil}\hfill\quad #3\par 
}

% Makro, ktere vyrobi cely obsah.
% Pouzit na zacatku, pred prvnim prispevkem.
% (Slo by upravit i pro pouziti na konci.)

\def\Contents{%
  \newpage
%  \immediate\closeout\obs
  \centerline {\LARGE {\bf Index of Talk Titles}}%Large
  \bigskip \bigskip
  \bgroup
    \parindent=-\parindent
    \spaceskip=.3333em
    \rightskip=5em plus 1fil
    \leftskip=-\parindent plus 0fil
    \parfillskip=-5em
    \input \jobname.obs
  \egroup
  \bigskip\bigskip
  \immediate\openout\obs=\jobname.obs
}

\def\ContentsSpeakers{%
  \newpage
%  \immediate\closeout\spk
  \centerline {\LARGE {\bf Index of Speakers}}%Large
  \bigskip \bigskip
  \bgroup
    \parindent=-\parindent
    \spaceskip=.3333em
    \rightskip=5em plus 1fil
    \leftskip=-\parindent plus 0fil
    \parfillskip=-5em
    \input \jobname.spk
  \egroup
  \bigskip\bigskip
  \immediate\openout\spk=\jobname.spk
}

% Makro pro zalomeni radku v nadpisu (\\ z nejakeho duvodu nechodi).

\def\CR{\penalty-10000}

% Promenna pro pomocny soubor.

\newwrite\obs
\newwrite\spk  % speakers
%\immediate\openout\obs=\jobname.obs


\Contents

%%%%%%%%%%%%%%%%%%%%%%%%%%%%%%%%%%%%%%%%%%%%%%%%%%%%%%%%%%%%%%%%%%%%

\ContentsSpeakers

%%%%%%%%%%%%%%%%%%%%%%%%%%%%%%%%%%%%%%%%%%%%%%%%%%%%%%%%%%%%%%%%%%%%
\newpage

%%%%%%%%%% Without abstract!
%\speaker{Ariana Angjeli}{}{Oakland University}
%\title{Linearly many faults in dual-cube-like networks}
%\endtitle
%The dual cube were introduced as better interconnections network
%than the hypercubes for large scale distributed memory multiprocessors. In
%this talk we introduce a generalization of these networks, called
%dual-cube-like networks, which preserve the basic structure of dual cubes
%and retain many of its topological properties. We investigate structural
%properties of these networks beyond simple measures such as connectivity. We
%prove that if up to $kn-(k(k+1))/2$ vertices are deleted from dual-cube-like
%n-regular network then the resulting graph will either be connected or will
%have large components and small components having at most $k-1$ vertices in
%total,and this result is sharp for $k\leq n$. As an application we derive
%additional results such as the cyclic vertex-connectivity and the restricted
%vertex-connectivity of these networks.
%
%Joint work with EDDIE CHENG AND LASZLO LIPTAK.


\speaker{Béla	Bollobás}{a}{The University of Memphis and Cambridge University}
\title{Bootstrap Percolation -- Problems, Results and Conjectures}
\endtitle

Bootstrap percolation is one of the simplest cellular automata around:
it can be viewed as a very simple model of the spread of an infection,
and is also closely related to the Ising model. In the past three decades,
much work has been done on bootstrap percolation on finite grids of a
fixed dimension, in which the initially infected set $A$ is obtained
by selecting its vertices at random, with the same probability $p$,
independently of all other choices. The focus has been on the {\em critical
probability} $p_c$, the value of $p$ at which the probability of {\em percolation}
(eventual full infection)  is $1/2$, say.

\vspace{10pt}
\noindent
Although bootstrap percolation used to be the hunting ground of statistical
physicists and probabilists only, it really belongs to probabilistic combinatorics.
My aim in this talk is to review of some of the basic results concerning
critical probabilities due to Aizenman, Lebowitz, Schonman, Cerf, Cirillo,
Manzo, Holroyd and others, and say a few words about the more recent results
proved by Balogh, Morris, Duminil-Copin and myself. I also hope to mention
some very recent results I have obtained with Holmgren, Smith, Uzzell and Balister
on the time a random initial set takes to percolate. In addition to the `proper'
results, I shall mention several lightweight problems and difficult conjectures.


%
%==================
%
%\speaker{Deepak Bal}{Carnegie Mellon University}
%\title{Packing Tree Factors in Random and Pseudo-Random Graphs}
%\endtitle
%For a fixed graph $H$ with $t$ vertices, an $H$-factor of a graph
%$G$ with $n$ vertices is a collection of vertex disjoint (not necessarily
%induced) copies of $H$ in $G$ covering all vertices of $G$.  We prove that
%certain pseudo-random and random graphs may have almost all of their edges
%covered by a collection of edge disjoint $H$-factors in the case where $H$
%is a tree.
%
%Joint work with Alan Frieze, Michael Krivelevich and Po-Shen Loh. 
%
%=================
%
%\speaker{Arthur Busch}{ University of Dayton}
%\title{Toughness and Hamiltonicity in k-trees}
%\endtitle
%A chordal graph is a $k$-tree if every maximal clique has order
%$k+1$ and every minimal separating set has order $k$.  In this talk we
%improve a theorem of Broesma, Xiong \& Yoshimoto on toughness in $k$-trees
%and consider some related corollaries. 
%
%
%==================
%
%\speaker{Steve Butler}{ Iowa State University}
%\title{Throttling zero forcing propagation speed}
%\endtitle
%Zero forcing is a game played on a graph that starts with a
%coloring of the vertices as white and black and at each step any vertex
%colored black with a unique neighbor colored white ``forces'' the color of
%the white vertex to be come black.  In this talk we look at what happens
%when we balance the size of the initial set of vertices colored black and
%the length of time that it takes for all vertices to be colored black.  We
%also give an example that shows it is possible in some graphs to slow down
%the speed of propagation in the graph by choosing larger initial sets.
%Joint work with Michael Young.
%
%
%==================
%
%\speaker{Ye Chen}{ West Virginia University}
%\title{r-hued Coloring of K4-minor Free Graphs}
%\endtitle
%A list assignment $L$ of $G$ is a mapping that assigns every
%vertex $v\in V(G)$ a set $L(v)$ of positive integers. For a given list
%assignment $L$ of $G$, an ($L,r$)-coloring of $G$ is a proper coloring $c$
%such that for any vertex $v$ with degree $d(v)$, $c(v)\in L(v)$ and $v$ is
%adjacent to at least min$\{d(v),r\}$ different colors. The {\it $r$-hued
%chromatic number} of $G$, $\chi_{r}(G)$, is the least integer $k$ such that
%for any $v \in V(G)$ with $L(v) = \{1,2,\cdots,k\}$, $G$ has an
%$(L,r)$-coloring. The {\it $r$-hued list chromatic number} of $G$,
%$\chi_{L,r}(G)$, is the least integer $k$ such that for any $v \in V(G)$ and
%every list assignment $L$ with $|L(v)| = k$, $G$ has an $(L,r)$-coloring.
%Let $K(r)=r+3$ if $2 \le r \le 3$, and  $K(r)=\lfloor 3r/2\rfloor+1$ if
%$r\ge 4$. We proved that if $G$ is a $K_4$-minor free graph, then
%
%(i) $\chi_r(G) \le K(r)$, and the bound can be attained;
%
%(ii) $\chi_{L,r}(G) \le K(r)+1$.
%
%
%===============
%
%\speaker{Zhihong Chen}{ Butler University}
%\title{Eulerian-connected Graphs}
%\endtitle
%A graph $G$ is Eulerian-connected if for any $u$ and $v$ in
%$V(G)$, $G$ has a spanning $(u,v)$-trail.
%A graph $G$ is edge-Eulerian-connected if for any $e'$ and $e''$ in $E(G)$,
%$G$ has a spanning $(e',e'')$-trail.
%For an integer $r\ge 0$, a graph is called $r$-Eulerian-connected if for any
%$X\subseteq E(G)$ with $|X|\le r$,
%and for any $u,\ v \in V(G)$, 
%$G$ has a spanning $(u,v)$-trail $T$ such that $X\subseteq E(T)$. The
%$r$-edge-Eulerian-connectivity of a graph can be defined similarly.
% Let $\theta (r)$ be the minimum value of $k$ such
%that every $k$-edge-connected graph is $r$-Eulerian-connected.
% Catlin proved that $\theta(0)=4$. In this talk, we shall show that
%$\theta(r)=4$ for $0\le r\le 2$, and $\theta(r)=r+1$ for $r\ge 3$, and show
% some new results on $r$-edge-Eulerian connectivity. 
%
%===============
%
%\speaker{Ilkyoo Choi}{ University of Illinois at Urbana-Champaign}
%\title{On Choosability with Separation of Planar Graphs with Forbidden}
%\endtitle
%We study choosability with separation which is a constrained
%version of list coloring of graphs. A \emph{$(k,d)$-list assignment} $L$ of
%a graph $G$ is a function that assigns to each vertex $v$ a list $L(v)$ of
%at least $k$ colors and for any adjacent pair $xy$, the lists $L(x)$ and
%$L(y)$ share at most $d$ colors. A graph $G$ is $(k,d)$-choosable if there
%exists an $L$-coloring of $G$ for every $(k,d)$-list assignment $L$. This
%concept is also known as choosability with separation. We prove that planar
%graphs without 4-cycles are $(3,1)$-choosable and that planar graphs without
%5-cycles and 6-cycles are $(3,1)$-choosable. In addition, we give an
%alternative and slightly stronger proof that triangle-free planar graphs are
%$(3,1)$-choosable. 
%
%
%===============
%
%\speaker{John Engbers}{ University of Notre Dame}
%\title{ Extremal $H$-colorings of forests and trees}
%\endtitle
%Given a finite graph $H$, an $H$-coloring of a finite, simple
%graph $G$ (or graph homomorphism) is a map from the vertices of $G$ to the
%vertices of $H$ that preserves edge adjacency.  $H$-colorings generalize
%many important graph theoretic notions, such as proper $q$-colorings (via $H
%= K_{q}$) and independent sets (via $H$ as an edge with a loop on one
%endvertex).
%
%Given a family of graphs $\mathcal{G}$ and a graph $H$, which graph(s) in
%$\mathcal{G}$ have the largest number of $H$-colorings?  We present several
%results for the family of $n$-vertex forests and the family of $n$-vertex
%trees.  Numerous open questions remain.
%
%
%================
%
%\speaker{David Galvin}{ University of Notre Dame}
%\title{Twin conventions and graph Stirling numbers}
%\endtitle
%Recently Griffiths asked (and answered) the question: in how many
%ways can $n$ sets of twins at a twin convention break into non-empty groups,
%with no group allowed to contain a pair of twins? One approach to this
%problem leads naturally to the notion of a ``graph Stirling number of the
%second kind''. I'll say what we know about these numbers and what we don't,
%and then use them to address Griffith's question with ``twins'' replaced by
%``triplets'', ``quadruplets'', etc. Partly joint work with Do Trong Thanh.
%
%===============
%
%\speaker{Sogol Jahanbekam}{ University of Illinois at Urbana-Champaign}
%\title{Rainbow Spanning Subgraphs in Edge-Colored Complete Graphs}
%\endtitle
%For integers $n$ and $t$, let $r(n,t)$ be the maximum number of
%colors in an
%edge-coloring of the complete graph $K_n$ that does not contain $t$
%edge-disjoint rainbow spanning trees.  Let $s(n,t)$ be the maximum number of
%colors in an edge-coloring of $K_n$ containing no rainbow spanning subgraph
%with diameter at most $t$.  We determine r(n,t) for
%$n>2t+\sqrt{6t-\frac{23}{4}}+\frac52$ and for $n=2t$. We also determine
%$s(n,t)$ for all integers $n$ and $t$. 
%This is joint work with Douglas B. West. 
%
%
%
%===============
%
%\speaker{Zeinab Maleki}{ University of Illinois at Urbana-Champaign}
%\title{Some lower bounds for the intersection dimension of graphs}
%\endtitle
%The {\it intersection dimension} of a graph with respect to a set
%of non-negative integers~$L$ is the small number~$l$ for which there is an
%assignment on the vertices to subsets $A_v \subseteq \{1,\dots, l\}$, such
%that every two vertices $u,v$ are adjacent if and only if $|A_u \cap A_v|\in
%L$.
%The {\it bipartite intersection dimension} is defined similarly when the
%conditions are considered only for the vertices in different parts. The {\it
%absolut dimension} of a graph $G$ is the minimum intersection dimension of
%$G$ over all sets $L$. Finding graphs with large (bipartite) absolute
%dimension would have important consequences in the complexity theory.
%In this talk, we present some lower bounds for the (bipartite) intersection
%dimension of a graph with respect to various types $L$ in terms of the
%minimum rank of graph.
%(This is a joint work with Behnaz Omoomi)
%
%
%
%
%===============
%
%
%\speaker{ Terry McKee}{Wright State University}
%\title{ Weighty Characterizations of Two Graph Classes}
%\endtitle
%Defining the weight of a complete subgraph to be the number of
%maxcliques that contain it, I give a stylish characterization of strongly
%chordal graphs in terms of their clique graphs, and a related-but-quirky
%characterization of trivially perfect graphs. 
%
%
%====================
%
%\speaker{ Darren Parker}{Grand Valley State University}
%\title{ Multidesigns for Graph-Triples of Order 6 }
%\endtitle
%We call $T=(G_1,G_2,G_3)$ a \emph{graph-triple of order t} if the
%$G_i$ are pairwise non-isomorphic graphs on $t$ non-isolated vertices whose
%edges can be combined to form $K_t$.  If $m \geq t$, we say $T$
%\emph{divides} $K_m$ if $E(K_m)$ can be partitioned into copies of the
%graphs in $T$ with each $G_i$ used at least once, and we call such a
%partition a \emph{$T$-multidecomposition}.  In this talk, we determine
%$T$-multidecompositions of complete graphs, where $T$ is a graph-triple of
%order 6. Moreover, we determine maximum multipackings and minimum
%multicoverings when $K_m$ does not admit a multidecomposition. 
%
%====================
%
%\speaker{ Nicholas Peterson}{The Ohio State University}
%\title{  On Random $k$-Out Graphs with Preferential Attachment }
%\endtitle
%We generalize a model of Hansen and Jaworski for random mappings
%which exhibit preferential attachment to a mapping $[n]\mapsto[n]^k$ - or,
%equivalently, a random digraph $M_{n,k}^{\alpha}$ with labeled arcs and
%uniform out-degree $k$. Each vertex starts with some weight $\alpha>0$; each
%vertex chooses $k$ images, one at a time in a fixed order, with probability
%of choosing a given vertex proportional to its current weight; and the
%weight of the chosen vertex increases by 1 before the next decision is made.
%The limiting case $M_{n,k}^{\infty}$ is a uniformly random $k$-out digraph
%with labeled arcs.
%
%We establish a limiting distribution for the vertex connectivity of the
%graph obtained from $M_{n,k}^{\alpha}$ by ignoring loops, multiple edges,
%and arc directions. We also seek to answer the question: how fast must
%$\alpha$ grow (relative to $n$) in order to make the difference between
%$M_{n,k}^{\alpha}$ and $M_{n,k}^{\infty}$ asymptotically negligible?
%Measuring with the total variation distance, we establish
%$\alpha=\Theta(\sqrt{n})$ as a sharp threshold for this behavior.
%
%This is a joint work with Boris Pittel. 
%
%====================
%
%\speaker{ Gregory Puleo}{University of Illinois at Urbana-Champaign}
%\title{  Tuza's Conjecture for Graphs of Max Average Degree Less than 7 }
%\endtitle
%Suppose I wish to make a graph $G$ triangle-free by removing a
%small
%number of edges.  An obvious obstruction is the presence of a large
%set of edge-disjoint triangles, since I must remove one edge from each
%triangle. On the other hand, removing all the edges in a maximal set
%of edge-disjoint triangles clearly makes $G$ triangle-free.
%\emph{Tuza's conjecture} states that the worst-case number of edges
%that must be removed is somewhere between these extremes: if $\tau(G)$
%is the number of edges that must be removed to make $G$ triangle-free
%and $\nu(G)$ is the maximum number of edge-disjoint triangles in $G$,
%then Tuza's conjecture states that $\tau(G) \leq 2\nu(G)$. Using the
%method of discharging, we show that Tuza's conjecture holds whenever
%the maximum average degree $\mathrm{Mad}(G) < 7$. This subsumes
%several earlier results and represents the first application of
%discharging to this problem. 
%
%
%
%================
%
%\speaker{ Michael Santana}{University of Illinois at Urbana-Champaign}
%\title{  Pairs of forbidden subgraphs for pancyclicity}
%\endtitle
%The Matthew-Sumner Conjecture claims that every 4-connected,
%claw-free graph contains a hamiltonian cycle.  This conjecture has provided
%the impetus for a great deal of research into structural properties of
%claw-free graphs, and in particular, pancyclicity.  While it is known that
%forbidding the claw is an insufficient condition to imply pancyclicity of
%4-connected graphs, we look to forbid more substructures to obtain this end.
%In this talk we will present a characterization of the pairs of forbidden
%subgraphs that imply pancyclicity for 4-connected graphs.  This extends a
%result of Gould, {\L}uczak, and Pfender for 3-connected graphs.  (Joint work
%with James Carraher, Michael Ferrara, and Timothy Morris) 
%
%====================
%
%\speaker{ Benjamin Stucky}{University of Oklahoma}
%\title{  An Algorithmic Approach to Raising the Lower Bound of $R(5,5)$ }
%\endtitle
%The primary focus of this talk is to explore some basic methods
%for pushing the lower bound on the classical Ramsey number $R(5,5)$ from
%$43$ to $44$.  Interestingly enough, we can phrase the problem in the
%following way: Is it possible to have an arrangement of $43$ people in a
%room such that no $5$ are \textit{either} mutual acquaintances \textit{or}
%mutual strangers (a question that perhaps gave rise to the informal
%description of this problem as the ``Party Problem'').  First, I explore the
%simple brute force method involving testing $2^{903}$ configurations of
%$2$-colorings of the edges of $K_{43}$, the complete graph on $43$ vertices.
%Next I discuss some simplifications to this approach, including attempting
%to exploit symmetry in configurations to eliminate graphs to test.  Finally,
%I discuss the approach involving extending an existing solution on one less
%vertex, drastically eliminating the number of graphs to be searched to
%$2^{42}$. Having been unsuccessful in using my algorithm to raise the lower
%bound of $R(5,5)$, I discuss some rationalizations for my ``failure'' and an
%argument for $R(5,5)=43$. I will start with a description and history of the
%problem, continue with a description of my approach, and finish by outlining
%my results. 
%
%
%
%
%====================
%
%\speaker{ Keke Wang}{West Viriginia University}
%\title{ Cycle chains and Hamiltonian 3-connected claw-free graphs}
%\endtitle
%We develop a cycle chain method to prove that every
%3-edge-connected graph G is supereulerian if every 3-edge cut of G
%intersects with short cycles in G. This is applied to the study of
%Hamiltonian claw-free graphs, and provide a unified treatment with short
%proofs for several former results. New sufficient conditions for Hamiltonian
%claw-free graphs are also obtained. 
%
%====================
%
%
%
%\speaker{ Minquan Zhan}{Millersville University}
%\title{ The Discharging Method and 3-Connected Essentially 10-Connected Line Graphs}
%\endtitle
%We use the discharging method to prove that every 3-connected,
%essentially 10-connected line graph is hamiltonian connected. 
%
%
%
%====================
%
%
%\speaker{ Meng Zhang}{West Virginia University}
%\title{ Supereulerian Graphs and The Petersen Graph}
%\endtitle
%A graph G is supereulerian if G has a spanning eulerian subgraph.
%Boesch et al proposed the problem of characterizing supereulerian graphs. In
%this talk I will present any 3-edge-connected graph with at most 11 edgecuts
%of size 3 is supereulerian if and only if it cannot be contractible to the
%Petersen graph. 
%
%====================
%
%
%\speaker{ Xiangqian "Joe" Zhou}{Wright State University}
%\title{ Unavoidable minors of large 4-connected bicircular matroids }
%\endtitle
%The bicircular matroid of a graph G is the matroid with ground
%set E(G) and a subset is independent iff each connected component of the
%subgraph spanned by that set contains at most one cycle.   In this talk, we
%present a recent result on unavoidable minors of large 4-connected
%bicircular matroids. This is joint work with Daniel Slilaty, Deborah Chun,
%and Tyler Moss.  
%
%
%=================
%
%\speaker{ Daniel Poole}{The Ohio State University}
%\title{ Weak Hamilton Cycles in Random Hypergraphs}
%\endtitle
%We say that a hypergraph, $H$, has a weak Hamilton cycle if there is some
%cyclic ordering of the vertices of $H$, such that each consecutive pair of
%vertices is contained in some hyperedge. We find the sharp threshold for the
%existence of a weak Hamilton cycle in the random d-uniform hypergraph,
%$H_d(n,m)$, on $n$ vertices with $m$ hyperedges. While the
%Erd\H{o}s-R\'{e}nyi graph $G(n,m)=H_2(n,m)$ with high probability develops a
%Hamilton cycle when the minimum vertex degree reaches $2$, for $d>2$
%$H_d(n,m)$ whp becomes Hamiltonian sooner, when the minimum vertex degree
%reaches $1$. Our proofs use hypergraph analogues of Pos\'{a}'s lemma and De
%la Vega's Theorem. 
%
%=================
%
%
%\speaker{ Tom Mahoney}{University of Illinois at Urbana-Champaign}
%\title{ Online sum list coloring of graphs}
%\endtitle
%In \textit{online list coloring} (introduced by Zhu and by Schauz in 2009), on each round the set of vertices having a particular color in their lists is revealed, and the coloring algorithm chooses an independent subset to receive that color. A graph $G$ is said to be $f$-paintable for a function $f: V(G) \to \mathbb{N}$ if there is an algorithm to produce a successful coloring when each vertex $v$ is allowed to be presented at most $f(v)$ times.
%
%In 2002 Isaak introduced \textit{sum list coloring} and the resulting parameter called \textit{sum-choosability}. The \textit{online sum-choosability}, or \textit{sum-paintability}, of $G$  is the least $\sum f(v)$ over all functions $f$ such that $G$ is $f$-paintable; this value is denoted by $\chi_{sp}(G)$.
%
%The \textit{generalied theta-graph} $\Theta_{\ell_1,\dots,\ell_k}$ consists of two vertices joined by internally disjoint paths of lengths $\ell_1,\dots,\ell_k$.
%Strengthening results of Berliner et al., we show that the sum-paintability of $G$ depends only on the sum-paintability of its blocks, and we prove $\chi_{sp}(K_{2,r})=2r+\min\{ l + m : lm>r \}$.
%The \textit{book} $B_r$ is the graph $\Theta_{1,2,\dots,2}$ with $r$ internally disjoint paths of length 2. We prove $\chi_{sp}(B_r)=2r+\min\{ l + m : m(l-m)+\binom{m}{2} > r \}$.
%Using these results, we determine the sum-paintability for all generalized theta-graphs.
%
%=================
%
%
%
%\speaker{ Axel Brandt}{University of Colorado Denver}
%\title{ 21st Century Cops and Robbers}
%\endtitle
%Since it's introduction in 1976, pursuit-evasion games on graphs appear in many forms. The well known cops and robbers game in which a `Cop' attempts to locate a `Robber' on a graph relates to significant theoretical topics such as treewidth.
%
%We consider a variation introduced by Seager in 2012. Each turn, the Cop `probes' a vertex $v$ in the pursuit graph $G$ and receives the Robber's current distance from $v$. The Robber is then permitted to move to any vertex adjacent to his current location, with the exception of $v$. The Cop wins if at any time he is able to determine the Robber's location exactly.
%
%Seager gave a non-constructive proof that on any tree $T$, the Cop can win in at most $|T|-2$ probes. In this talk, we give an explicit strategy for the Cop to win on a tree $T$ in at most $|T|-2$ probes, meeting Seager's bound.
%
%This is joint work with Jennifer Diemunsch, Catherine Erbes, Jordan Legrand, and Casey Moffatt.
%
%=================
%
%
%
%\speaker{ Jaehoon Kim}{University of Illinois at Urbana-Champaign}
%\title{ (0,1)-improper coloring of sparse triangle-free graph}
%\endtitle
%A graph $G$ is a $(0,1)$-colorable if $V(G)$ can be partitioned into two sets $V_0$ and $V_1$ so that $G[V_0]$ is an independent set and $G[V_1]$ has maximum degree at most $1$. The problem of verifying whether a graph is $(0,1)$-colorable is NP-complete even in the class of planar graphs of girth 9.
%
%Maximum average degree, $Mad(G)= \max_{H\subset G}\{\frac{2|E(H)|}{|V(H)|}\}$, is a graph parameter measuring how sparse the  graph $G$ is. Borodin and Kostochka showed that every graph $G$ with $Mad(G)\leq \frac{12}{5}$ is $(0,1)$-colorable, thus every planar graph with girth at least $12$ also is $(0,1)$-colorable.
%
%The aim of this talk is to prove that every triangle-free graph $G$ with $Mad(G)\leq \frac{22}{9}$ is $(0,1)$-colorable. We prove the slightly
%stronger statement that every triangle-free graph $G$ with $|E(H)|<\frac{11|V(H)|+5}{9}$ for every subgraph $H$ is $(0,1)$-colorable and show that there are infinitely many not  $(0,1)$-colorable graphs $G$ with $|E(G)|=\frac{11|V(G)|+5}{9}$. This is joint work with A. V. Kostochka and Xuding Zhu.
%
%
%=================
%
%
%
%\speaker{ Stephen Young}{University of Louisville}
%\title{ An Alon-Boppana result for the normalized Laplacian}
%\endtitle
%We prove a Alon-Boppana style bound for the spectral gap of the normalized
%Laplacian for general graphs via a bound on the weighted spectral radius of
%the universal cover graph.
%
%=================
%
%
%\speaker{ Robert Seiver}{Miami University}
%\title{ Set Families without Chain-traces}
%\endtitle
%Given a family $\cF$ of subsets of a set $X$ and $W\subseteq X$, the {\it trace} of $\cF$ on $W$, denoted by $\cF|_W$ is defined to be  $\cF|_W=\{F\cap W: F\in \cF\}$. If for some $k$-set $W$ in $X$,
%$\cF|_W$ contains a maximal chain $C_0\subset C_1\cdots \subset C_k$ with $|C_i|=i$,
%then we say that $\cF$ has a {\it $k$-chain-trace}. For fixed $k$ and sufficiently large $n$, Patk\'os  showed that every family $\cF$ of $k$-subsets of $[n]$ with more than $\binom{n-1}{k-1}$ members
%has a $k$-chain-trace. The bound is best possible. Patk\'os made a more general conjecture
%that for fixed $m\geq k\geq 2$ and sufficiently large $n$, every family $\cF$ of $m$-subsets of $[n]$ with
%more than $\binom{n-m+k-1}{k-1}$ members has a $k$-chain-trace. We prove Patk\'os' conjecture and for large $n$ also establish strong stability property of the unique
%extremal family.
%
%=================
%
%
%\speaker{ Huseyin Acan}{The Ohio State University}
%\title{ Evolution of a Random Permutation Graph}
%\endtitle
%Associated with a permutation $\sigma$ of $[n]$, there is a graph $G_{\sigma}$ whose vertex set is $[n]$ and whose edges correspond to the inversions of $\sigma$. Let $\sigma(n,m)$ denote a permutation chosen uniformly at random from all permutations of $[n]$ with $m$ inversions.  We find a growth process of a random permutation in which we obtain $\sigma(n,m)$ after the $m$-th step. We show that the connectedness probability of $G_{\sigma(n,m)}$ is non-decreasing in $m$. We also find the threshold value of the number of edges to be $\frac{6}{\pi^2}n\ln(n)$ for the connectedness of the inversion graph $G_{\sigma(n,m)}$. This is a joint work with Boris Pittel.
%
%=================
%
%
%
%\speaker{ Daniel McDonald}{University of Illinois at Urbana-Champaign}
%\title{ On-line rank number of trees}
%\endtitle
%A $k$-ranking of a graph $G$ is a labeling of its vertices from $[k]$ such that any nontrivial path whose endpoints have the same label contains a larger label.  The least $k$ for which $G$ has a $k$-ranking is the rank number of $G$, also known as tree depth.  Applications of rankings include VLSI design, parallel computing, and factory scheduling. The on-line ranking problem asks for an algorithm for ranking the vertices of $G$ as they are presented one at a time along with all previously ranked vertices and the edges between (so each vertex is presented as the lone unranked vertex in a partially labeled induced subgraph of $G$ whose final placement in $G$ is not specified).  The on-line rank number of $G$ is the minimum over all such algorithms of the largest label that algorithm can be forced to use.  We give upper and lower bounds on the on-line rank number of trees in terms of maximum degree, diameter, and other structural parameters.
%
%
%=================
%
%
%
%\speaker{ Peter Blanchard}{Grinnell College, University of Iowa}
%\title{ Coloring in Classical Lattices}
%\endtitle
%We prove a Van der Waerden type theorem concerning finite colorings of the Boolean Lattice $B_n$ and the Young Lattice.  This is joint project with Amitava Bhattacharya, Tata Institute of Fundamental Research, Mumbai, India.
%
%
%=================
%
%
%\speaker{ Lisa Espig}{Carnegie Mellon University}
%\title{ Threshold for Zebraic Hamilton Cycles in Random Graphs}
%\endtitle
%When studying random graphs, we often want to know the threshold for the emergence of a particular structure. In other words, in the Erd\"os-R\'enyi random graph model $G_{n,p}$ - where we consider a graph on $n$ vertices with each edge appearing randomly with probability $p$, how large must $p$ be in order to guarantee a matching in the graph with high probability? This and other thresholds are well-studied. Here we find the threshold for a different kind of structure. Namely, we will find the threshold for which a randomly 2-colored random graph contains a zebraic Hamilton cycle - one whose edges alternate between the two colors.
%
%
%=================
%
%
%\speaker{ Bernard Lidicky}{University of Illinois at Urbana-Champaign}
%\title{ Coloring planar graphs with 4-triangles}
%\endtitle
%A sharpening of Gr\" otzsch Theorem, the Gr\" unbaum-Axenov
%Theorem, states that
%every planar graph with at most three triangles is 3-colorable.
%It is best possible since not all planar graphs with four triangles
%are 3-colorable.
%In this talk, we discuss 3-colorability of planar graphs
%with four triangles.
%
%This is joint work with O. V. Borodin, Z Dvo\v{r}\'ak, A. Kostochka, and M. Yancey
%
%
%=================
%
%
%\speaker{ Uta Ziegler}{Western Kentucky University}
%\title{ Embedding of 4-regular planar graphs on a grid}
%\endtitle
%The author and her collaborators have developed methods to embed 4-regular planar graphs into a 2-dimensional grid while preserving the topology (i.e the cyclic order of vertices at each vertex). We call a 4-regular planar graph algebraic if it can be reduced to a trivial graph (consisting of one 4-regular vertex) by the repeated collapses of digons. Here we present a special case of our algorithm for algebraic 4-regular graphs, where it can be shown that the length of a constructed embedding is linear in the number of vertices of the original graph. These graph embedding results can be applied to obtain results about the rope length of knots.


\newpage

\centerline {\LARGE {\bf List of Participants}}
\participant{Ghodratollah}{Aalipour}{alipour.ghodratollah@gmail.com}{Sharif University of Technology and UTEP}
\participant{Naman}{Agarwal}{naman33k@gmail.com}{UIUC}
\participant{László}{Babai}{laci@cs.uchicago.edu}{University of Chicago}
\participant{Paul N.}{Balister}{pbalistr@memphis.edu}{University of Memphis}
\participant{József}{Balogh}{jobal@math.uiuc.edu}{UIUC}
\participant{Ali}{Belabbas}{belabbas@illinois.edu}{UIUC ECE}
\participant{Frank}{Bernhart}{farb45@gmail.com}{retired}
\participant{Scott}{Binski}{smbinski@memphis.edu}{University of Memphis}
\participant{Csaba}{Biro}{csaba.biro@louisville.edu}{University of Louisville}
\participant{Derek}{Boeckner}{boecknerd@gmail.com}{University of Nebraska-Lincoln}
\participant{Béla}{Bollobás}{b.bollobas@dpmms.cam.ac.uk}{University of Memphis and Cambridge University}
\participant{Steve}{Butler}{butler@iastate.edu}{Iowa State University}
\participant{Jane}{Butterfield}{butter@umn.edu}{University of Minnesota}
\participant{James}{Carraher}{s-jcarrah1@math.unl.edu}{University of Nebraska--Lincoln}
\participant{Guantao}{Chen}{gchen@gsu.edu}{Georgia State University}
\participant{Ilkyoo}{Choi}{ichoi4@illinois.edu}{University of Illinois at Urbana-Champaign}
\participant{Jonathan}{Cutler}{cutlerjo@mail.montclair.edu}{Montclair State University}
\participant{Shagnik}{Das}{shagnik@ucla.edu}{UCLA}
\participant{Louis}{DeBiasio}{debiasld@miamioh.edu}{Miami University}
\participant{Michelle}{Delcourt}{delcour2@illinois.edu}{University of Illinois, Urbana-Champaign}
\participant{Dwight}{Duffus}{dwight@mathcs.emory.edu}{Emory University}
\participant{Zdeněk}{Dvořák}{ook@ucw.cz}{Georgia Tech and Charles University in Prague}
\participant{Aly}{El Gamal}{alyelgamal@gmail.com}{ECE, UIUC}
\participant{John}{Engbers}{jengbers@nd.edu}{University of Notre Dame}
\participant{Asaf}{Ferber}{ferberasaf@gmail.com}{Tel Aviv University}
\participant{Zoltán}{Füredi}{z-furedi@math.uiuc.edu}{Rényi Mathematical Institute}
\participant{David}{Galvin}{dgalvin1@nd.edu}{University of Notre Dame}
\participant{Wenying}{Gan}{wgan@math.ucla.edu}{UCLA}
\participant{John}{Goldwasser}{jgoldwas@math.wvu.edu}{West Virginia University}
\participant{Andrey}{Grinshpun}{agrinshp@mit.edu}{MIT}
\participant{Andrzej}{Grzesik}{Andrzej.Grzesik@uj.edu.pl}{Jagiellonian University}
\participant{Esfandiar}{Haghverdi}{ehaghver@indiana.edu}{Indiana University Bloomington}
\participant{Tigran}{Hakobyan}{hakobya2@illinois.edu}{University of Illinois at Urbana-Champaign}
\participant{Jie}{Han}{jhan22@gsu.edu}{Georgia State University}
\participant{David}{Hannasch}{David.A.Hannasch@gmail.com}{UIUC}
\participant{Stephen}{Hartke}{hartke@math.unl.edu}{University of Nebraska-Lincoln}
\participant{Penny}{Haxell}{pehaxell@math.uwaterloo.ca}{University of Waterloo}
\participant{Anil}{Hirani}{hirani@illinois.edu}{University of Illinois}
\participant{Ping}{Hu}{pinghu1@illinois.edu}{University of Illinois at Urbana-Champaign}
\participant{Sogol}{Jahanbekam}{jahanbe1@illinois.edu}{University of Illinois at Urbana-Champaign}
\participant{Tao}{Jiang}{jiangt@miamioh.edu }{Miami University}
\participant{Adam}{Jobson}{asjobs01@louisville.edu}{University of Louisville}
\participant{Richard}{Johnson}{rabjohnson@gmail.com}{University of Memphis}
\participant{Tomas}{Juskevicius}{tomas.juskevicius@gmail.com}{University of Memphis}
\participant{Hemanshu}{Kaul}{kaul@iit.edu}{Illinois Institute of Technology}
\participant{Mitch}{Keller}{mitch.keller@unl.edu}{University of Nebraska-Lincoln}
\participant{Lauren}{Keough}{s-lkeough1@math.unl.edu}{University of Nebraska - Lincoln}
\participant{Hal}{Kierstead}{kierstead@asu.edu}{Arizona State University}
\participant{Jaehoon}{Kim}{kim805@illinois.edu}{UIUC}
\participant{Ringi}{Kim}{ringikim@princeton.edu}{Princeton University}
\participant{Teeradej}{Kittipassorn}{t.kittipassorn@memphis.edu}{University of Memphis}
\participant{Alexandr}{Kostochka}{kostochk@math.uiuc.edu}{University of Illinois at Urbana-Champaign}
\participant{Lucas}{Kramer}{ljkramer@iastate.edu}{Iowa State University}
\participant{Daniel}{Král'}{d.kral@warwick.ac.uk}{University of Warwick}
\participant{Andre}{Kundgen}{akundgen@csusm.edu}{California State University San Marcos}
\participant{Choongbum}{Lee}{cb_lee@math.mit.edu}{Massachusetts Institute of Technology}
\participant{Jonathan}{Lee}{jdl43@cam.ac.uk}{DPMMS, University of Cambridge}
\participant{John}{Lenz}{lenz@math.uic.edu}{University of Illinois at Chicago}
\participant{Linda}{Lesniak}{lindalesniak@gmail.com}{Drew University/Western Michigan University}
\participant{Shoham}{Letzter}{sl597@cam.ac.uk}{University of Cambridge}
\participant{Nana}{Li}{nli3@gsu.edu}{Georgia State University}
\participant{Bernard}{Lidický}{lidicky@illinois.edu}{University of Illinois}
\participant{Jonathan}{Ligo}{ligo2@illinois.edu}{University of Illinois at Urbana-Champaign}
\participant{Hong}{Liu}{hliu36@illinois.edu}{UIUC}
\participant{Allan Siu Lun}{Lo}{s.a.lo@bham.ac.uk}{Birmingham University}
\participant{Sarah}{Loeb}{sloeb2@illinois.edu}{UIUC}
\participant{Jie}{Ma}{jiema@math.ucla.edu}{UCLA}
\participant{Siva Theja}{Maguluri}{siva.theja@gmail.com}{UIUC}
\participant{Thomas}{Mahoney}{tmahone2@illinois.edu}{UIUC}
\participant{Zeinab}{Maleki}{zmaleki@ilinois.edu}{University of Illinois Urbana-Champaign}
\participant{Ryan}{Martin}{rymartin@iastate.edu}{Iowa State University}
\participant{Andrew}{McConvey}{mcconve2@illinois.edu}{University of Illinois at Urbana-Champaign}
\participant{Daniel}{McDonald}{dmcdona4@illinois.edu}{UIUC}
\participant{John}{McSorley}{jmcsorley@math.siu.edu}{Southern Illinois University Carbondale}
\participant{Kevin}{Milans}{milans@math.wvu.edu}{West Virginia University}
\participant{Zevi}{Miller}{millerz@miamioh.edu}{Miami University}
\participant{Theodore}{Molla}{tmolla@asu.edu}{Arizona State University}
\participant{Timothy}{Morris}{timothy.morris@ucdenver.edu}{University of Colorado Denver}
\participant{Kevin}{Moss}{kmoss@iastate.edu}{Iowa State University}
\participant{Dhruv}{Mubayi}{mubayi@math.uic.edu }{University of Illinois at Chicago}
\participant{Richard}{Mycroft}{r.mycroft@bham.ac.uk}{Birmingham University}
\participant{Humberto}{Naves}{hnaves@math.ucla.edu}{UCLA}
\participant{Andrew}{Newman}{newmanj2@miamioh.edu}{Miami University}
\participant{Vladimir}{Nikiforov}{vnikifrv@memphis.edu}{University of Memphis}
\participant{Vivek}{Nittoor}{vivek@is.s.u-tokyo.ac.jp}{The University Of Tokyo}
\participant{Sergey}{Norin}{snorin@math.mcgill.ca}{McGill University}
\participant{Steven}{Osborne}{sosborne@iastate.edu}{Iowa State University}
\participant{Cory}{Palmer}{ctpalmer@illinois.edu}{University of Illinois, Urbana-Champaign}
\participant{Anna}{Payne}{apayne1@ashland.edu}{Ashland University}
\participant{Michael}{Pelsmajer}{pelsmajer@iit.edu}{Illinois Institute of Technology}
\participant{Bhargav}{Peruvemba Narayanan}{bp338@cam.ac.uk}{University of Cambridge}
\participant{Šárka}{Petříčková}{sarpetka@gmail.com}{UIUC}
\participant{Florian}{Pfender}{Florian.Pfender@ucdenver.edu}{UC Denver}
\participant{Mike}{Plantholt}{mikep@ilstu.edu}{Illinois State University}
\participant{Noah}{Prince}{nprince@imsa.edu}{Illinois Mathematics and Science Academy}
\participant{Landon}{Rabern}{landon.rabern@gmail.com}{Arizona State University}
\participant{Jamie}{Radcliffe}{aradcliffe1@math.unl.edu}{UNL}
\participant{Alexander}{Razborov}{razborov@cs.uchicago.edu}{University of Chicago}
\participant{Steven}{Reich}{stevenreich47@gmail.com}{University of Pittsburgh}
\participant{Ben}{Reiniger}{reinige1@illinois.edu}{UIUC}
\participant{Bruce}{Reznick}{reznick@math.uiuc.edu}{UIUC}
\participant{Andrzej}{Ruciński}{rucinski@amu.edu.pl}{Emory University and Adam Mickiewicz University}
\participant{Wojciech}{Samotij}{wojteksa@gmail.com}{Cambridge UK and Tel Aviv University}
\participant{Michael}{Santana}{santana@illinois.edu}{University of Illinois at Urbana-Champaign}
\participant{Andrew}{Schwartz}{aschwartz@semo.edu}{Southeast Missouri State University}
\participant{Maryam}{Sharifzadeh}{sharifz2@illinois.edu}{UIUC}
\participant{Gaurav}{Sinha}{gsinha@caltech.edu}{California Institute of Technology}
\participant{Papa}{Sissokho}{psissok@ilstu.edu}{Illinois State University}
\participant{David}{Smith}{dav@ugcs.caltech.edu}{Grand Valley State University}
\participant{Richard}{Snyder}{rjs84@pitt.edu}{University of Pittsburgh}
\participant{Hannah }{Spinoza}{hkolb2@illinois.edu}{University of Illinois}
\participant{Christopher}{Stocker}{christopher.stocker@marquette.edu}{Marquette University}
\participant{Derrick}{Stolee}{stolee@illinois.edu}{University of Illinois}
\participant{Benny}{Sudakov}{bsudakov@math.ucla.edu }{University of California (LA)}
\participant{Gábor N.}{Sárközy}{gsarkozy@wpi.edu}{Worcester Polytechnic Institute}
\participant{Michael}{Tait}{mtait@math.ucsd.edu}{University of California - San Diego}
\participant{Craig}{Timmons}{ctimmons@ucsd.edu}{University of California San Diego}
\participant{Charles}{Tomlinson}{ctomlinson2@math.unl.edu}{University of Nebraska at Lincoln}
\participant{William T.}{Trotter}{trotter@math.gatech.edu}{Georgia Institute of Technology}
\participant{Konstantinos}{Tyros}{ktyros@math.toronto.edu}{Department of Mathematics, University of Toronto}
\participant{Ali}{Vakilian}{vakilia2@illinois.edu}{UIUC}
\participant{Jennifer}{Vandenbussche}{jvandenb@spsu.edu}{Southern Polytechnic State University}
\participant{Jacques}{Verstraëte}{jbaverstraete@gmail.com}{University of California (SD)}
\participant{Jan}{Volec}{j.volec@warwick.ac.uk}{University of Warwick}
\participant{Dominik}{Vu}{dominik.vu@memphis.edu}{University of Memphis}
\participant{Ruidong}{Wang}{rwang49@math.gatech.edu}{Georgia Institute of Technology}
\participant{Paul}{Wenger}{pswsma@rit.edu}{Rochester Institute of Technology}
\participant{Douglas}{West}{west@math.uiuc.edu}{Zhejiang Normal University and University of Illinois}
\participant{Jeffrey}{Wheeler}{jwheeler@pitt.edu}{University of Pittsburgh}
\participant{Jennifer}{Wise}{jiwise2@illinois.edu}{University of Illinois Urbana-Champaign}
\participant{Hehui}{Wu}{noshellwhh@gmail.com}{McGill University}
\participant{Yihong}{Wu}{yihongwu@illinois.edu}{University of Illinois Urbana-Champaign}
\participant{Jiaming}{Xu}{jxu18@illinois.edu}{University of Illinois at Urbana-Champaign}
\participant{Rupei}{Xu}{xuxxx852@umn.edu}{University of Minnesota--Twin Cities}
\participant{Derrek}{Yager}{yagerdjd@gmail.com}{UIUC}
\participant{Matthew}{Yancey}{yancey1@illinois.edu}{UIUC}
\participant{Amy}{Yates}{ayates2@gsu.edu}{Georgia State University}
\participant{Elyse}{Yeager}{yeager2@illinois.edu}{UIUC}
\participant{Liana}{Yepremyan}{liana.yepremyan@mail.mcgill.ca}{McGill University}
\participant{Stephen}{Young}{stephen.young@louisville.edu}{University of Louisville}
\participant{Gexin}{Yu}{gyu@wm.edu}{College of William and Mary}
\participant{Cheng}{Zhao}{cheng.zhao@indstate.edu}{Indiana State University}
\participant{Yufei}{Zhao}{yufeiz@mit.edu}{MIT}
\participant{Lidong}{Zhou}{lidongz@andrew.cmu.edu}{Carnegie Mellon University}
\participant{Xuding}{Zhu}{xudingzhu@gmail.com}{Zhejiang Normal University}


\end{document}

%%%%%%%%%%%%%%%%%%%%%%%%%%%%%%%%%%%%%%%%%%%%%%%%%%%%%%%%%%%%%%%%%%%%
